\documentclass[12pt]{article}
\usepackage{amsmath,amsfonts,fullpage,amsthm,graphicx}

\newtheorem*{theorem}{Theorem}
\newtheorem*{fact}{Fact}
\newtheorem*{goal}{Goal}
\newtheorem*{idea}{Idea}
\newtheorem*{defn}{Definition}
\newtheorem*{ex}{Example}

%Style section
%\setlength{\textheight}{10in}
%\setlength{\textwidth}{7in}
%\hoffset=-0.75in
 \pagestyle{plain}
% \setlength{\topmargin}{0in}
% \setlength{\headsep}{0in}
% \setlength{\oddsidemargin}{0in}
% \setlength{\evensidemargin}{0in}

\begin{document}


\pagestyle{empty} %use this to get rid of page numbers



\noindent
{Math 166 \hskip 1 truein  Name:\ \hrulefill}

\noindent
%\vskip 0.3truein
%\bigskip
%\noindent
{An additional study of series}
\noindent \hspace{2.5in} Section Number: \hrulefill
\bigskip

\begin{goal} \rm
Explore even more methods of determining whether series converge or diverge. 
\end{goal}
%\smallskip

\begin{enumerate}
%\setcounter{enumi}{-1}
\item Name all the tests for convergence/divergence of series in our series toolbox. Put a * by the tests that only work for series with positive terms.\\ Note, your answer should be longer than last worksheet!
%\item %A \textbf{sequence} is a~ \rule{.8in}{.4pt} ~of numbers.  A sequence converges if its~ \rule{.8in}{.4pt} ~have a limit.
%\item A \textbf{series} is a~ \rule{.6in}{.4pt} ~of numbers. A series converges if its~ \rule{.8in}{.4pt} \rule{.2in}{.4pt} \rule{.8in}{.4pt} %\rule{.5in}{.4pt} ~has a limit.
%\item
%\end{enumerate}
\vspace{1in}

\item Use any of the tests you listed in 1.\ to determine whether the following series converge or diverge. Be sure to verify the hypotheses, state your conclusion, and say which test you used. If you can, list a second test that might work as well.
\begin{enumerate}
\item $\displaystyle\sum_{n=1}^\infty \frac{(-1)^n}{4^{2n+1}(n+1)}$

\vfill


\item $\displaystyle\sum_{n=1}^{\infty} \frac{e^n}{n!}$

\vfill
%\hrule
\item $\displaystyle\sum_{n=3}^\infty \frac{(-12)^n}{n}$

\vfill
\newpage

\item $\displaystyle\sum_{n=1}^{\infty} \frac{5^n}{n^n}$ 

(Hint: When you see something to the $n$th power in a series, it's often a good idea to use the root test.)

\vfill


\item $\displaystyle\sum_{n=1}^\infty \left(2^{-n}+\frac{3n+5}{n^2}\right)$


%\newpage


%\vfill

%\item $\displaystyle\sum_{n=2}^\infty \frac{1}{n(\mathrm{ln}(n))^2}$

%\newpage

\vfill

\item $\displaystyle\sum_{n=2}^\infty \frac{n}{(n^3-1)^{\frac{3}{7}}}$
\vfill

\end{enumerate} 

\end{enumerate}

\end{document}
